\chapter{Building a thought experiment \label{chapter:thoughtexp}}

As discussed in the introduction\cmt{``As discussed in the introduction'' - not yet}, any comparison between a \gls{ge} and \gls{abm} approach to location theory needs to explain what it aims to achieve from doing so. One stated aim, for example, is to find out what \gls{abm} is actually good at. Does it have any strengths that make it able to analyse problems beyond the reach of an analytic approach? In short, is \gls{abm} better at some things than \gls{ge}, and vice versa?

But this begs the question: what does `better' look like? The job of this section is to find some candidate answers to that question.

%%%%%%%%%%%%%%%%%%%%%%%%%%%%%%%%%%%%
% WHAT ARE ABSTRACT MODELS GOOD FOR?
%%%%%%%%%%%%%%%%%%%%%%%%%%%%%%%%%%%%
\section{What are abstract models good for? \label{sec:abstractmodels}}

Abstract models allow for concentrating on the relationships between variables in the model: to come to a full understanding of how they work, separate from calibrating against real-world data, which can actually cloud the ability to see what is going on, if data is not used in the correct way.

There is no point in simply adding data to a model. The problem here is that it can always be done: data can be added in a myriad of ways. But there are only a very few ways for it to be used appropriately, and the tendency to mistake `includes data' for `is therefore more realistic' is strong.
%I wonder if the Bali model would have been accepted without actual rainfall data? The abstract picture: less appealling to people who want to see simulacra of reality. 

Of course, ultimately, reality is the arbiter of all models. The purpose of this chapter, however, is to ask exactly how this comparison to reality should inform the model builder, and what role the `thought experiment' model has.

My conclusion's going to be something like: the commonly-held view of model `validation' goes back to statistical tests against data: `fitting' is commonly held to be a suitable goal. However, the reality of model-building, interpretation and use in the social sciences is much less clear-cut. The role of the modellers' interpretation plays a much larger role, as does the form of `praxis' involved in particular disciplines. Missing the fact of this praxis may be responsible for why critics often mistake economists' models for their actual understanding of the world, rather than tools that aid that understanding.
\cmt{`tools that aid understanding': repeated that line here. Make sure to delete it from one of the two places it currently resides.}

%%%%%%%%%%%%%%%%%%%%%%%%%%%%%%%%%%%
% DOES PARSIMONY = POWERFUL?
%%%%%%%%%%%%%%%%%%%%%%%%%%%%%%%%%%%
\section{Does `more parsimonious' equal `more powerful?' \label{parsimony}}

A model, by definition, is simpler than the system it aims to represent. \fixref{slightly worried I read this point somewhere, which makes it plagiarism via amnesia. have a think, see if it's traceable.} Throughout the previous chapters, a more complex reality has been described, as well as justifications for limiting models to... blah.

%was in space cost change section...
So what has just been argued? It could be this: space is complex - far too complex, in fact, to be amenable to the kind of abstract modelling this thesis proposes as feasible. Is this a fair assessment? Perhaps \gls{ge} can offer some lessons here. The tradition of economic modelling that \gls{ge} carries on is comfortable with its own simplifications, while openly discussing what compromises are made in abstracting away from a more complex reality. \gls{abm} has tended instead towards more complexity in a kind of theoretical mission creep. Edmonds, for example, wants to drop Occam's Razor because `simplicity is not truth-indicative' (\fixref{edmonds p.1})  - and he sees nothing intrinsically valuable about avoiding complexity:

\begin{quote}
``... model selection `for the sake of simplicity' is either: simply laziness; is really due to pragmatic reasons such as cost or the limitations of the modeller; or is really a relabelling of more sound reasons due to special circumstance or limited data. Thus appeals to it should be recognised as either spurious, dishonest or unclear and hence be abandoned.'' (p.8)
\end{quote}
%my job here, over the thesis, is to demonstrate that simplicity is necessary to help address the question at hand: how should we model the impact of space cost changes on the `general' system?



%%%%%%%%%%%%%%%%%%%%%%%%%%%%%%%%%%%
% CONSTRUCT A LANGUAGE, SUBSTANTIVE PROPOSITIONS
%%%%%%%%%%%%%%%%%%%%%%%%%%%%%%%%%%%
\section{`Construct a language that will yield as many substantive propositions as possible.' \label{sec:constructalanguage}}

\fixref{Citing from http://www.mbs.murdoch.edu.au/workingpapers/163.html - need proper refs, and go read Friedman myself.}

Milton Friedman offers one answer. As Leeson notes, Friedman argues that theories are tautologous forms of language - they cannot directly provide their own criteria for success: ``the relevant question to be asked is usefulness and not rightness or wrongness'' (Friedman 1953 p.7 quoted in Leeson.) They must function as `engines of analysis', producing propositions that \emph{can} be held up to the world. This leads Friedman to say that one should - 

\begin{quote}
``Construct a language that will yield as many substantive propositions as possible.'' (Friedman 1962 p.8)
\end{quote}

Might this be a good candidate for answering this section's question - what does a `better' theory look like? Clearly, it is not possible to ask \gls{ge} and \gls{abm} to churn out `substantive propositions' and find out which has the most - but this does offer a tool for comparison. What propositions can each give? Can they both say something about propositions the other makes, or are they producing very different ones?

The \emph{kind} of propositions any particular 'engine of analysis' can produce matters too. \gls{ge} comes from an international trade perspective - it wants to be able to say something about why places trade, and why they can end up with such disparate wealth. A key proposition it gives is this: in the model, it is possible for regions to diverge due to endogenous forces; there is no need for there to be any difference in initial factor endowment. \cmt{mentioning GE's `substantive propositions' here - worth cross-reffing with wherever I end up writing them up in detail.} Notice that sentence started with `in the model'; the next question must be, `is it possible to find some element of reality that can offer support to this proposition, or falsify it?'

\cmt{What I want to argue somewhere: current thought on transition and energy is offering a whole bag of propositions. What I want theory to do is to give us substantive propositions that confirm or deny the statements made, or offer interesting new ones that nevertheless might shed light on those issues. GE deals with a very different proposition 'base'. It's all about trying to explain the core-periphery pattern as 'an endogenous process...' Is it really any good at answering the kind of questions I'm interested in?}

%nice summary of the f-twist here: http://motls.blogspot.com/2009/02/milton-friedmans-f-twist.html
%This will need writing up much better and more carefully, if I do end up using it. I think I probably should.
Friedman's chief point here is one that has caused a great deal of consternation, and offered an easy line of attack for theorists unhappy with the simplifying assumptions of economics. Called the `F-twist' by Samuelson (see \citealt{blaug_methodology_1992} pp.91 for an overview), the caricature version says: `model assumptions don't matter, and don't have to be realistic. The only thing that matters is predictive power.' This is an unfair interpretation. Einstein, for example, made the same point as regards physics:

\begin{quote}
``Physics constitute a logical system of thought which is in a state of evolution, whose basis cannot be distilled, as it were, from experience by an inductive method, but can only be arrived at by free invention. The justification (truth content) of the system rests in the verification of the derived propositions by sense experiences. The skeptic will say: `it may well be true that this system of equations is reasonable from a logical standpoint. But it does not prove that it corresponds to nature'. You are right, dear skeptic. Experience alone can decide on truth.'' (A. Einstein, Ideas and Opinions, New York, 1960, pp. 322 and 355 (quoted by Kornai, op. cit., pp. 9-10), quoted in \citealt{kaldor_irrelevance_1972} p.1239.)
\end{quote}

However, there are very separate challenges involved for what this means in physics and in economics. Whether or not the universe `really is' as relativity describes has quite different consequences to whether people `really are' rational, constrained optimisers. Some of these are beyond the scope of this thesis: for example, that the conception of human nature a theory uses may feed back into policy structures that end up creating what it describes (see \fixref{jc scott map thing} scott...).

To illustrate the `social reality' problem, consider an example. Imagine a model of allotment production versus growing food in one's backyard: the factors can be reduced to time and distance. Does an actor stand to gain more utility from travelling to the allotment and putting time in there, or by staying at home and growing, thus eliminating travel costs? (Note that this is quite similar to the choice facing actors in the backyard models in section \ref{sec:backyard}.) For a set of actors who can choose between allotment and home-garden plots of similar size, rationally, it seems they should choose to eliminate travel costs and produce in their own backyard. But what if reality doesn't conform to this - there is a set of people who choose allotments over their own, perfectly suitable garden. Now one might propose an alteration to the theory, making sure to keep the assumption of rationality: utility is gained from leisure use of the home garden that would be lost by turning it into a vegetable patch. That would need to be outweighed by the disutility of distance before an actor would dig up their begonias.

But is that proposition true? The risk is this: whatever outcome actually exists in the world, the theorist \emph{starts} from the assumption that it must be the product of rational actors. The job of the model is then to conform to the world. This is, of course, tautologous: rational decisions lead to the outcomes seen in reality, and reality confirms that, indeed, rational decisions are the cause. No analytic insight results; indeed, it seems to actually nullify the point of making the model in the first place, and imposes a false 'model' on a reality that may well be far more involved. What role might social convention play, for example?

This suggests that, whatever role 'substantive propositions' from 'model engines' might play, it is not entirely straightforward. \cmt{I think I'm heading back towards the 'dialogue between model and world' stuff here. Though it's also that the model might predict well. And whether utility is a good way of getting how human system choices might work. Cf. example of pool player modelling. Where did I read that example?}

To return to another point of Friedman's: there is often confusion about what 'predictive power' means. This is dealt with in the next section.

%%%%%%%%%%%%%%%%%%%%%%%%%%%%%%%%%%%
% SHOULD THEORIES PREDICT?
%%%%%%%%%%%%%%%%%%%%%%%%%%%%%%%%%%%
%ontological prediction etc
\section{Should theories predict? \label{prediction}}

On the face of it ,the `thought experiment' form of model dealt with in this thesis - and most common in economics - cannot meet the most obvious test: is it better at prediction? This is not the case, however. Two quick definitions will serve to illustrate this. Prediction is usually taken to mean \emph{future} prediction. This can instead be defined as \emph{forecasting}: making a claim about something that will occur at a future point. (These sort of predictions are often, but not always, probabalistic.) Betz suggests another very useful term: ontological prediction (\citealt{betz_prediction_2006}.) He gives the example of the prediction of Neptune's existence from Newton's laws - Uranus' orbit implied that another body must exist. An ontological prediction, according to Betz, is `timeless' - the phenomenon was always there. Einstein's predictions about light bending near the sun, for example, described something that always happened, it just wasn't looked for - and doubtless Eddington wouldn't have thought \emph{how} to look, without the theory.

This fits Epstein's distinction (\citealt{epstein_why_2008}) between earthquakes and tectonic plates. Discovery of the latter was as a result of the success of an ontological prediction. Note also, this ontological prediction does indeed open the way to forecasting. It places a boundary around possible future outcomes that were unknown before: earthquakes will mostly happen on or near lines of tectonic plates. This also leads to the concept of epicentres, which can be used to model tsunamis.

J Stephen Lansing and James Kremer's model of Balinese rice management is a good example of an ontological prediction from a very simple simulation \cite{lansing_priests_2007}. Attempts to apply Green Revolution-style growing techniques to Bali were actually damaging the existing rice-growing system. Lansing had tried to tell the authorities about the role that the Balinese water temples played in finding an optimal trade-off between fallow and growing that maximised yield and minimised pests, to no avail. However, when Kremer helped Lansing model the system, and the key dynamic became easy to explain, they had more success. This was an ontological claim: a simple model that demonstrates a mechanism that manages rice production, in a way that appears to capture the salient features of the real-world system. 

This is, perhaps, what Epstein claims is happening in `generative' social science.\cmt{`This is, perhaps, what Epstein claims is happening in `generative' social science.' This is its first mention.} If so, Epstein is wrong to claim that `growing it' can count as `explaining it'.  Generating a model phenomenon offers only a hypothesis - one of Friedman's `substantive propositions'. The tautology of the model still needs to be theoretically tied to the real world.  Does this mean models are \emph{only} analytic engines that output hypthoses, not explanatory tools? 

%%%%%%%%%%%%%%%%%%%%%%%%%%%%%%%%%%%
% THEORY WITHOUT MEASUREMENT?
%%%%%%%%%%%%%%%%%%%%%%%%%%%%%%%%%%%
\section{Theory without measurement? \label{Koopman}}

Koopman�s scathing article attacks what he sees as a work that is `unbendingly empiricist in outlook' (\citealt{koop
 p.172): collecting data and looking for trends - in this case, business cycles - with no thought given to underlying mechanism. (Indeed, Burns and Mitchell, who he criticises, overtly eschewed theory [cf. review, follow up])

Koopman compares the progress of economic analysis to that of studying planetary motion, from Tycho Brahe's exhaustive but theoryless collection of data, through Kepler's ideas about circular motion informing his empirical work, finally to Newton's formulation of `fundamental laws'. The discovery of `empirical regularities' he calls the `Kepler stage', and of `fundamental laws' the `Newton stage.' The one is distinguished from the other by being `at once more elementary and more general.'

Despite seemingly approving of the need for a `Kepler stage', Koopman politely attacks the empirical void in Burns and Mitchells' analysis of business cycles; they aim for:

\begin{quote}
``� faithful observation and summarizing of the cyclical characteristics of a large number of economic series. The toolkit of the theoretical economist is deliberately spurned. Not a single demand or supply schedule or other equation expressing the behavior of men or the technical laws of production is employed explicitly in the book, and the cases of implicit use are few and far between.'' (ibid p.163)
\end{quote}

They might find statistical rules for predicting future business cycles, they might not; Koopmans still doubts the validity of the whole enterprise. Koopmans would, presumably, share the queaziness with which geographical economists view spatial interaction models. Clearly, such models effectively show that an `empirical regularity' exists, yet they propose no mechanism - they haven't, in Koopman�s terms, progressed to the Newton stage.

...

What, then, is the role of mechanism in models? `Kepler stage' meteorological models attempted to base forecasts on statistical inference, and were not at all successful. It was only one mechanisms were proposed, and built into models, that things started to improve. \fixref{refs for meteorology - stats to mechanics?} Mechanism, then, at the very least appears to offer a predictability premium for certain systems. 

But the objection of neoclassical economists to Kepler-stage analyses doesn�t appear to flow from any concern with predictability. So where does this uncompromising dedication to microfoundations and mechanism come from? Can similar reasons be found in earlier and more recent economic analyses? Can the answer help shed any light on what kind of model choices should be made?

\cmt{What I need to be heading towards here: setting the discussion of mechanism in the context of the representative agent problem.}

%%%%%%%%%%%%%%%%%%%%%%%%%%%%%%%%%%%%%%%
% THE DOUGHNUT OF EMPIRICAL CORRECTNESS
%%%%%%%%%%%%%%%%%%%%%%%%%%%%%%%%%%%%%%%
\section{The doughnut of empirical correctness \label{doughnut}}

``The doughnut of empirical correctness in a theory constitutes its worth, while its hole of untruth constitutes its weakness.'' Samuelson in Blaug 1980 p.97.

%%%%%%%%%%%%%%%%%%%%%%%%%%%%%%%%%%%%%%%
% Maths: models, reality, choices
%%%%%%%%%%%%%%%%%%%%%%%%%%%%%%%%%%%%%%%
\section{Using the maths in the models: what does it mean? \label{sec:mathsinmodels}}

\cmt{What I want to do here: e.g. run through examples where an increasing returns equation may appear to be rubbish, and explain how one might decide how to use it. So e.g. love of variety: actors can increase their utility infinitely by splitting their production into infinite slices of different goods. Monocomp rules that out, but since we can use LoV to make space work, that condition is not necessary. Also, LoV and 1d vs 2d space...}

\subsection{Problems with `love of variety' \label{sec:LoVprobs}}

If an actor producing in their backyard has a `love of variety' described using a CES function, they can make themselves as happy as they like by the simple expedient of splitting their day and producing a different good in each timeslot. 



%%%%%%%%%%%%%%%%%%%%%%%
% CUTTINZ
%%%%%%%%%%%%%%%%%%%%%%%
\begin{comment}
One might wonder: what good are substantive propositions if they turn out to be nonsense? Surely the point is that they need to be tested against reality. Indeed, this is true. The aim here is to argue for a clear blue line between the 'engine of analysis', the propositions it can churn out, and job of comparing them to the world. \cmt{What the bejeesus am I talking about here?} 
...
Bali stuff:

If one were cynical, there's an argument to say all this shows is that policymakers have an unhealthy propensity to prefer "quantitative-looking answers", regardless of their actual philosophical merit. This may be self-serving, but I prefer to think that "ontological models" like this *can* reveal hidden phenomena in exactly the way Betz describes, and that - as the Bali example shows - this revealing of dynamics previously invisible to the eyes of policymakers is an entirely sound justification for modelling.

(Cf. the paper above - the point isn't the simplicity of the model, but the ability of the modeller to theoretically justify its connection to real-world features - what Craik calls a "similar relation structure.") The model was convincing enough to get official support for traditional methods of rice management. 

There's also another very important reason to model, I'd argue, due to the double-edged nature of models: they can do exactly the opposite of the Bali model - keep things invisible by insisting they've already captured everything important about a system. Friedman's classic positivist argument underlines this: it doesn't matter at all, he tells us, whether the assumptions of one's model are realistic or not - as long as the model predicts successfully. That entirely fails in economics, where powerful policy bodies have consistently argued e.g. that a country liberalised too fast, or too slow - but at any rate, it certainly wasn't the fault of the model they were using if things didn't work out.  Actually, Friedman is right - but as regards ontological prediction. Uncovering previously unseen dynamics in social systems - rather than forecasting - is a very useful outcome of a model.  

%%%%%%%%%%%
% Friedman bits
%%%%%%%%%%%
Blaug on Walras: asks whether he was wedded to formalism at the expense of realism. �How seriously we take this charge of formalism all depends of course on what we think Walras was doing.� He concludes most analysts see Walras as - through the editions of his Elements, �an increasing tendency of Walras to fit the world to the model rather than the model to the world.� 

(note Schumpeter�s praise of Walras, and Hayek�s attack on Schumpeter below.)

Milton Friedman, in contrast - ``his problem is the problem of form, not of content: of displaying an idealised picture of the economic system, not of constructing an engine for analysing concrete problems.'' [quoted in ibid p.569, though Blaug gives no reference, annoyingly. Hmm� seems to be from (Friedman 1955, Leon Walras and his System: A Review Article. American Economic Review XLV.1: 900-909.]

\end{comment}

